% !TEX root = ../main.tex
\section{Conclusion}

This paper formulates multi-UAV wildfire rescue as a closed-loop decision
problem connecting aerial situation understanding, resource-constrained
planning, and runtime adaptation. Based on this formulation, we develop a
modular system that integrates Vision Analysis, Mission Planning, Validation,
Episodic and Lesson Memory, and Runtime Replanning. We also construct a
real-UAV-informed AirSim evaluation suite containing 80 fixed episodes across
four difficulty levels and six complementary mission metrics.

Across four downstream foundation-model configurations, the complete framework
maintains high basic task completion, while execution quality and mission
continuity decline as operational pressure increases. Controlled ablations
show that Validation provides an occasional reliability safeguard, Memory
becomes increasingly valuable as resource and scheduling pressure grows, and
Runtime Replanning is critical under dynamic perturbations. Fine-tuning
Qwen3-8B with validated planning demonstrations retains 98.24\% of the
Gemini~3.1~Pro planner's average score on simple and normal episodes, while
reducing mean planning time by approximately 63\% and planning-token usage by
73\%. These results support lower-cost local deployment of the initial planning
module.

The current evidence remains limited to a calibrated simulation, a fixed
episode set, and one real-platform configuration. Moreover, the specialized
planner has not yet been evaluated under hard and expert runtime conditions.
The evaluation also does not include a direct numerical comparison with an
external end-to-end rescue method because existing systems assume different
inputs, action spaces, resource constraints, and evaluation protocols. The
reported evidence therefore supports the framework within the task and
evaluation setting defined here, rather than establishing general superiority
over traditional planning methods. Future work will extend the evaluation to
real-flight and mixed real--sim trials, more diverse platforms and fire
dynamics, and larger fleets, while developing more selective memory retrieval
and more robust, efficient replanning for high-pressure missions.

\endinput
